\section{The Algorithmic Plan's Position in the Ensemble}
\label{sec:percentile}

\subsection{Percentile Definition}

Let $\mathcal{E} = \{P_1, P_2, \ldots, P_N\}$ be an ensemble of $N$ valid
redistricting plans for a state with $k$ districts.  Let $f: \mathcal{P} \to
\mathbb{R}$ be a scalar statistic (e.g., Polsby-Popper mean, Democratic seat
share, minority-majority district count).  The \emph{percentile} of the
algorithmic plan $P^*$ with respect to $f$ and $\mathcal{E}$ is:

\begin{equation}
  \pi_f(P^*) = \frac{1}{N} \#\{i \in [N] : f(P_i) \leq f(P^*)\}.
  \label{eq:percentile}
\end{equation}

A plan at the 50th percentile is typical; plans at the 5th or 95th percentile
are unusual but not necessarily extreme.  Courts have generally treated plans
above the 99th or below the 1st percentile as requiring explanation
\citep{herschlag2020quantifying}.

\subsection{Statistics of Interest}

The G-track considers three families of statistics:

\begin{enumerate}
  \item \textbf{Polsby-Popper (PP) compactness.}
    For a district with area $A$ and perimeter $P$,
    $\pp = 4\pi A / P^2 \in (0, 1]$.  Higher is more compact.
    We report the plan-level mean $\bar{\pp}$ and the minimum
    $\pp_{\min}$ across districts.

  \item \textbf{Partisan seat share.}
    The number of districts won by each party under a
    specified electoral baseline (e.g., 2020 presidential vote).
    We report Democratic seats $S_D \in \{0, 1, \ldots, k\}$.

  \item \textbf{Minority-majority districts (MM).}
    The count of districts with minority voting-age population exceeding
    50\%.  We report $\mathrm{MM} \in \{0, 1, \ldots, k\}$.
\end{enumerate}

\subsection{Interpreting Percentile Position}

The percentile framework supports two distinct claims:

\begin{keybox}
\textbf{Claim 1 (Compactness not extreme).}
If $\pi_{\bar\pp}(P^*) \in [0.60, 0.90]$, the \ar{} plan is more compact
than average but not an outlier.  This is \emph{legally desirable}: a plan
at the 99th percentile of compactness might be suspected of packing
communities to achieve partisan ends via the ``compactness route.''

\textbf{Claim 2 (Partisanship near median).}
If $\pi_{S_D}(P^*) \in [0.40, 0.60]$, the partisan outcome of the \ar{}
plan is indistinguishable from the geographic median.  This is the
strongest possible evidence that the plan is not gerrymandered.
\end{keybox}

We will see in G.1 and G.2 that \ar{} achieves both claims for the majority
of studied states, with notable exceptions in states with strong geographic
partisan sorting (Georgia, Texas).

\subsection{Sensitivity to Ensemble Choice}

The percentile $\pi_f(P^*)$ depends on the choice of ensemble $\mathcal{E}$.
Different chain types (flip vs.\ \recom{}) produce different distributions;
different ensemble sizes affect statistical precision; different population
tolerance bounds affect plan diversity.  In G.1, we use published ensemble
results from \citet{herschlag2020quantifying} for North Carolina and
\citet{deford2021recombination} for other states, acknowledging that these
are \recom{} ensembles with $N \geq 10{,}000$ and population tolerance
$\pm 0.5\%$.

Where published ensembles are unavailable, we use literature-reported
distributional summaries (mean, standard deviation, approximate percentile
bounds) to bound $\pi_f(P^*)$.  These bounds are marked with \est{} throughout
the G-track papers.
